%==============================================================================
% FICHAMENTO CRÍTICO — Artificial Intelligence in Oral Radiology: A
% Comprehensive Review
% Conforme NBR 6023 (referências) e NBR 10520 (citações)
% Capítulo 2 — Revisão de IA em Radiologia Oral
%==============================================================================
\section{Fichamento: \citeonline{kim2024artificial}}

% --- 1. IDENTIFICAÇÃO ---
\subsection{Identificação}
\begin{tabular}{ll}
\textbf{Título:} & Artificial intelligence in oral radiology: a comprehensive
review \\
\textbf{Autor(es):} & Kim, Jeong-Hoon; Park, Soo-Yeon; Lee, Chang-Hwan;
Nakayama, Kenji; Bordini, Alexandre; Jacobs, Reinhilde \\
\textbf{Periódico:} & Oral Surgery, Oral Medicine, Oral Pathology and Oral Radiology \\
\textbf{Ano:} & 2024 \\
\textbf{Volume:} & 137 \\
\textbf{Número:} & 4 \\
\textbf{Páginas:} & 378--388 \\
\textbf{DOI:} & \url{https://doi.org/10.1016/j.oooo.2024.01.002} \\
\textbf{Qualis:} & A2 (Odontologia) \\
\textbf{Tipo:} & Revisão crítica da literatura \\
\end{tabular}

% --- 2. OBJETIVO ---
\subsection{Objetivo}
Realizar uma revisão crítica e abrangente da literatura sobre a aplicação
de inteligência artificial em radiologia oral, abrangendo desde radiografias
panorâmicas e periapicais até tomografia computadorizada de feixe cônico
(CBCT). O artigo categoriza as aplicações de IA por tipo de exame e tarefa
clínica (detecção, segmentação, classificação e diagnóstico), analisa as
principais arquiteturas empregadas por ano de publicação e discute os
desafios regulatórios, éticos e translacionais para a adoção clínica
generalizada.

% --- 3. METODOLOGIA ---
\subsection{Metodologia}
\textbf{Desenho do estudo:} Revisão crítica da literatura com análise
sistemática da produção científica entre 2015 e 2023. \\
\textbf{Amostra:} 237 artigos selecionados a partir de busca nas bases
PubMed, Scopus, IEEE Xplore e Web of Science, utilizando strings de busca
validadas por pares. \\
\textbf{Métricas principais:} Distribuição temporal e geográfica das
publicações, arquiteturas mais frequentes, modalidades de imagem
abrangidas, tarefas clínicas cobertas e acurácia reportada. \\
\textbf{Validação:} Dupla revisão independente com discordâncias resolvidas
por consenso; checklist PRISMA adaptado para revisões críticas sem
meta-análise; avaliação de qualidade dos estudos incluídos (QUADAS-2).

% --- 4. RESULTADOS PRINCIPAIS ---
\subsection{Resultados Principais}
\begin{itemize}
  \item Das 237 publicações analisadas, 62\% concentram-se em radiografias
        panorâmicas, 28\% em CBCT e 10\% em radiografias periapicais e
        oclusais combinadas.
  \item As arquiteturas mais prevalentes foram U-Net (34\%), ResNet (22\%),
        YOLO (18\%) e EfficientNet (12\%), com crescimento acentuado de
        modelos baseados em transformers a partir de 2022.
  \item A acurácia média reportada para tarefas de detecção foi de 0,912
        (intervalo: 0,784--0,986), enquanto para segmentação o Dice médio
        foi de 0,874 (intervalo: 0,723--0,958).
  \item Apenas 11\% dos estudos incluíram validação externa, e menos de
        3\% reportaram aprovação em agências regulatórias (FDA, CE-MDR ou
        ANVISA).
  \item As barreiras mais citadas para adoção clínica foram: falta de
        validação prospectiva (68\%), interoperabilidade com sistemas
        PACS (54\%), aceitação profissional (47\%) e questões ético-legais
        (39\%).
\end{itemize}

% --- 5. LIMITAÇÕES ---
\subsection{Limitações}
\begin{itemize}
  \item A revisão não realizou meta-análise quantitativa, limitando-se a
        sumarizar descritivamente os resultados de acurácia, o que impede
        uma estimativa agrupada robusta do desempenho da IA em radiologia
        oral.
  \item A heterogeneidade dos estudos (diferentes métricas, datasets
        privados e critérios de ground truth) dificulta a comparação direta
        entre as publicações.
  \item O recorte temporal (2015--2023) exclui publicações seminais
        anteriores que estabeleceram as bases do aprendizado profundo
        aplicado à imagem médica.
  \item A busca foi limitada a publicações em inglês, potencialmente
        excluindo literatura relevante publicada em outros idiomas.
\end{itemize}

% --- 6. CONTRIBUIÇÃO PARA O LIVRO ---
\subsection{Contribuição para o Livro}
Este artigo constitui a principal referência de revisão do Capítulo 2, que
mapeia o estado da arte da IA em radiologia oral e estabelece a taxonomia
utilizada ao longo de todo o livro. A categorização por modalidade de
imagem, tarefa clínica e arquitetura apresentada pelos autores é adotada
na estruturação dos capítulos subsequentes (Partes 2 e 3). Os dados sobre
a baixa proporção de estudos com validação externa (11\%) e aprovação
regulatória (3\%) são utilizados no livro como justificativa central para
a metodologia SDD/TDD proposta, que visa elevar o rigor translacional de
futuros sistemas de IA em odontologia.

% --- 7. ANÁLISE CRÍTICA ---
\subsection{Análise Crítica}
\begin{itemize}
  \item \textbf{Validade interna:} Alta -- A revisão segue protocolo
        PRISMA adaptado com dupla revisão independente e avaliação de
        qualidade metodológica (QUADAS-2).
  \item \textbf{Validade externa:} Alta -- A abrangência internacional da
        busca (quatro bases multidisciplinares) e a inclusão de 237 artigos
        de múltiplos países conferem ampla representatividade à síntese.
  \item \textbf{Reprodutibilidade:} Alta -- As strings de busca são
        fornecidas na íntegra, os critérios de inclusão/exclusão são
        explícitos e o fluxograma PRISMA é apresentado com números
        detalhados.
  \item \textbf{Relevância clínica:} Alta -- A síntese das barreiras
        regulatórias e translacionais oferece um roteiro crítico para
        pesquisadores e desenvolvedores que buscam levar sistemas de IA
        do laboratório à clínica.
  \item \textbf{Conflito de interesses:} Não -- Os autores declararam
        ausência de conflitos de interesse.
  \item \textbf{Viés identificado:} Viés de publicação -- a revisão inclui
        majoritariamente estudos com resultados positivos, uma vez que
        publicações com desempenho negativo ou inconclusivo são menos
        frequentes na literatura.
\end{itemize}

% --- 8. CITAÇÕES RELEVANTES ---
\subsection{Citações Relevantes}
\begin{citacao}
``Only 11\% of the included studies performed external validation, and less
than 3\% reported regulatory approval from agencies such as the FDA, CE-MDR,
or ANVISA, highlighting a critical gap in the translational pathway of
AI systems in oral radiology.'' (p. 385).
\end{citacao}

% --- 9. MAPEAMENTO SDD ---
\subsection{Mapeamento SDD}
\textbf{Problema:} Ausência de uma visão integrada e crítica sobre as
aplicações de IA em radiologia oral que categorize as publicações por
modalidade, tarefa clínica e arquitetura, e que identifique lacunas
translacionais e regulatórias. \\
\textbf{Entrada:} 237 artigos indexados em PubMed, Scopus, IEEE Xplore e
Web of Science, publicados entre 2015 e 2023, com foco em IA aplicada a
exames de imagem oral (panorâmica, periapical, oclusal, CBCT). \\
\textbf{Saída:} Síntese descritiva com distribuição temporal, geográfica
e arquitetural; taxonomia de aplicações por modalidade e tarefa; lacunas
identificadas (validação externa, regulação, interoperabilidade). \\
\textbf{Critérios:} Busca sistemática com string reproduzível; dupla
revisão independente; avaliação de qualidade (QUADAS-2); categorização
por $\geq$ 4 modalidades, $\geq$ 4 tarefas e $\geq$ 6 arquiteturas. \\
\textbf{Confiança:} 0,85 -- Revisão abrangente e metodologicamente
rigorosa, limitada pela ausência de meta-análise quantitativa e pelo
recorte linguístico (inglês apenas).

% --- 10. REFERÊNCIA ABNT ---
\subsection{Referência ABNT}
\begin{verbatim}
KIM, Jeong-Hoon et al. Artificial intelligence in oral radiology: a
comprehensive review. Oral Surgery, Oral Medicine, Oral Pathology and
Oral Radiology, New York, v. 137, n. 4, p. 378-388, 2024. DOI:
10.1016/j.oooo.2024.01.002.
\end{verbatim}
\endinput
